% year: תשע"ו
% type: מבחן
% semester: ב
% moed: ג
% source: exams/מבחנים/תשעו/מועד ג תשעו.pdf
% number: 1
% points: 10
% categories: true-false, ivt

\begin{question}
\textbf{הוכיחו או הפריכו:} למשוואה הבאה יש פתרון ב-$[0,\pi]$:
\[ \ln(x+1)-\sin(x)=1 \]
\end{question}

\begin{hint}
הגדירו $g(x)=\ln(x+1)-\sin(x)-1$ ובדקו את סימנה בקצוות הקטע.
\end{hint}

\begin{solution}
\textbf{הטענה נכונה.} נגדיר $g(x)=\ln(x+1)-\sin x-1$. הפונקציה $g$ רציפה ב-$[0,\pi]$ כסכום של פונקציות רציפות ($\ln(x+1)$ רציפה עבור $x>-1$).

בקצוות:
\[ g(0)=\ln1-\sin0-1=-1<0,\qquad g(\pi)=\ln(\pi+1)-\sin\pi-1=\ln(\pi+1)-1>0, \]
כאשר האי-שוויון האחרון נובע מכך ש-$\pi+1>e$ (כי $\pi+1>4>e$), ולכן $\ln(\pi+1)>\ln e=1$. (מספרית $\ln(\pi+1)\approx1.421$.)

לפי משפט ערך הביניים (משפט בולצאנו) קיימת $c\in(0,\pi)$ שעבורה $g(c)=0$, כלומר $\ln(c+1)-\sin c=1$. לכן למשוואה יש פתרון בקטע.
\end{solution}
